\documentclass[10pt,a4paper]{amsart}
\usepackage[margin=19mm]{geometry}
\usepackage{amsmath,amssymb,mathtools}
\usepackage[scr=rsfs,cal=euler]{mathalfa}
\usepackage{xfrac}
\usepackage[unicode,hidelinks,bookmarks=false]{hyperref}
\newcommand{\ie}{i.e.\ }
\newcommand{\cf}{cf.\ }
\newcommand{\resp}{respectively\ }
\newcommand{\Subsection}{\subsection}
\providecommand{\nobreakdash}{\nobreak\hbox{-}}
\setlength{\emergencystretch}{2em}
\multlinegap0pt
\usepackage{bm}
\usepackage{bbm}
\usepackage{dsfont}
\usepackage{xspace}
\usepackage{cancel}
\usepackage{mathtools}
\usepackage{amsthm}
\makeatletter
\@ifundefined{thm}{\newtheorem{thm}{Theorem}}{}
\@ifundefined{lem}{\newtheorem{lem}[thm]{Lemma}}{}
\@ifundefined{prop}{\newtheorem{prop}[thm]{Proposition}}{}
\makeatother
\makeatletter
\newcommand{\Bcy}{\operatorname{B^{cy}}}
\newcommand{\Brep}{\operatorname{B^{rep}}}
\expandafter\ifx\csname Def\endcsname\relax
\newenvironment{Def}{\begin{defn}}%
\fi
\newcommand{\E}{\mathbb{E}}
\newcommand{\N}{\mathbb{N}}
\renewcommand{\P}{\mathbb{P}}
\providecommand{\Spec}[1]{\operatorname{Spec}(#1)} 	%
\newcommand{\Sphere}{\mathbb{S}}
\newcommand{\Striv}{\mathbb{S}^{\mathrm{triv}}}
\newcommand{\TC}{\mathrm{TC}}
\newcommand{\THH}{\mathrm{THH}}
\newcommand{\TP}{\mathrm{TP}}
\newcommand{\Z}{\mathbb{Z}}
\newcommand{\boxx}{\square}
\newcommand{\cC}{\mathcal{C}}
\newcommand{\cofib}{\operatorname{cofib}}
\newcommand{\colimit}{\mathop{\rm colim}}
\expandafter\ifx\csname ex\endcsname\relax
\newenvironment{ex}{\begin{exm}}%
\fi
\expandafter\ifx\csname exs\endcsname\relax
\newenvironment{exs}{\begin{exms}}%
{\hfill$\square$\end{exms}}
\fi
\newcommand{\horn}{\mathrm{horn}}
\newcommand{\infShv}{\mathcal{S}\mathrm{hv}}
\newcommand{\infSpc}{\mathcal{S}}
\newcommand{\infSpt}{\mathcal{S}\mathrm{pt}}
\newcommand{\lSch}{\mathbf{lSch}}
\newcommand{\lSm}{\mathbf{lSm}}
\expandafter\ifx\csname listabc\endcsname\relax
\newenvironment{listabc}{
                         \begin{list}{\alph{elno-abc})
                                     }{\usecounter{elno-abc}}
                      }{
                         \end{list}}
\fi
\newcommand{\logTC}{\mathrm{logTC}}
\newcommand{\logTHH}{\mathrm{logTHH}}
\newcommand{\logTP}{\mathrm{logTP}}
\expandafter\ifx\csname note\endcsname\relax
\newenvironment{note}{\begin{notat}}%
%{\hfill$\square$\end{notat}}

%\theoremstyle{definition}

%\renewcommand*\thechapter{\Roman{chapter}}
%\numberwithin{section}{chapter}
\fi
\expandafter\ifx\csname rem\endcsname\relax
\newenvironment{rem}{\begin{remk}}%
\fi
\expandafter\ifx\csname rems\endcsname\relax
\newenvironment{rems}{\begin{remks}}%
\fi
\expandafter\ifx\csname romanlist\endcsname\relax
\newenvironment{romanlist}{
                         \begin{list}{\roman{elno})
                                     }{\usecounter{elno}}
                      }{
                         \end{list}}
\fi
\newcommand{\tcofib}{\operatorname{tcofib}}
\newcommand{\tfib}{\operatorname{tfib}}
\makeatother
\title[Logarithmic motivic homotopy theory]{Logarithmic motivic homotopy theory}
\author{Federico Binda, Doosung Park, Paul Arne {\O}stv{\ae}r}
\date{Source: arXiv:2303.02729v3 (CC BY 4.0)}
\begin{document}
\maketitle

\noindent\emph{Source and licence.} Logarithmic motivic homotopy theory. Version arXiv:2303.02729v3 (CC BY 4.0), distributed under the Creative Commons Attribution 4.0 International licence, as recorded in the arXiv licence metadata for this exact version and shown on the abstract page https://arxiv.org/abs/2303.02729v3. Full rights notice: licenses/arxiv-2303.02729-CC-BY-4.0.txt.\par\smallskip

\noindent\textit{Editorial scope.} Counted unit: the Thm whose source label is \texttt{Hoch.13} in the article (source0.tex lines 10637--10647, source label \texttt{Hoch.13}), delivered together with the article's own complete original proof of it (source0.tex lines 10648--10835, one \texttt{proof} environment of 188 lines). No mathematics is written, repaired or re-derived: statement and proof are the article's own text line for line. Required context retained uncounted: logHprop.3 (lines 2089--2099); Hoch.37.1 (lines 9865--9878); Hoch.9 (lines 10007--10051); Hoch.2 (lines 10075--10082); Hoch.14 (lines 10123--10132); Hoch.12 (lines 10212--10223); Hoch.11 (lines 10595--10607); Thomdf.5.3 (lines 12816--12823); Thomdf.6 (lines 12915--12928); Thomdf.15 (lines 12958--12962). Every definition, display and label of those lines is retained unchanged. Omitted from this package and not redistributed: 1--2088, 2100--9864, 9879--10006, 10052--10074, 10083--10122, 10133--10211, 10224--10594, 10608--10636, 10836--12815, 12824--12914, 12929--12957, 12963--12999. Nothing inside a retained excerpt is omitted or altered: every retained body is delivered exactly as the article's lines are written, so no omission receipt of any kind accompanies this package. Citations: the retained excerpts carry no citation command at all, so no bibliography entry of the article is transcribed and no reference is redistributed. Reference adaptation: none was needed and none was made; every reference inside the delivered unit resolves inside it, so no unresolved reference or citation is produced. Printed slips kept literal: no printed word, symbol, hypothesis or number of the counted statement or of its proof was altered. Layout and class: the article's own class is \texttt{amsart} with packages including amsmath, amsfonts, amsthm, bbm, amssymb, standalone, comment, tikz-cd, xy, imakeidx, babel, enumerate, graphicx, stmaryrd; the wrapper is the lane's pinned \texttt{amsart} editorial wrapper at 10pt on A4, which restates the theorem-like environments the retained text needs and adds the article's own macro definitions that the retained text needs (\texttt{\textbackslash Bcy}, \texttt{\textbackslash Brep}, \texttt{\textbackslash E}, \texttt{\textbackslash N}, \texttt{\textbackslash P}, \texttt{\textbackslash Spec}, \texttt{\textbackslash Sphere}, \texttt{\textbackslash Striv}, \texttt{\textbackslash TC}, \texttt{\textbackslash THH}, \texttt{\textbackslash TP}, \texttt{\textbackslash Z}, \texttt{\textbackslash boxx}, \texttt{\textbackslash cC}, \texttt{\textbackslash cofib}, \texttt{\textbackslash colimit}, \texttt{\textbackslash horn}, \texttt{\textbackslash infShv}, \texttt{\textbackslash infSpc}, \texttt{\textbackslash infSpt}, \texttt{\textbackslash lSch}, \texttt{\textbackslash lSm}, \texttt{\textbackslash logTC}, \texttt{\textbackslash logTHH}, \texttt{\textbackslash logTP}, \texttt{\textbackslash tcofib}, \texttt{\textbackslash tfib}), with extra wrapper packages (bm, bbm, dsfont, xspace, cancel, mathtools). The delivered numbering is the unit's own. Assets: no retained span contains a figure environment, a graphics-inclusion command, a PDF-page inclusion, a table or a TikZ picture; no image file is copied into this package, the package declares no asset dependency and no image file is redistributed with it. Licence: version arXiv:2303.02729v3 of this article is distributed under the Creative Commons Attribution 4.0 International licence, as recorded in the arXiv licence metadata for this exact version and shown on the abstract page https://arxiv.org/abs/2303.02729v3. A full rights notice is delivered at licenses/arxiv-2303.02729-CC-BY-4.0.txt. Characters: every retained excerpt is pure ASCII, so the delivered engine, XeLaTeX with its default Latin Modern text font, reports no missing character.
\begin{prop}
\label{logHprop.3}
Suppose $S\in \lSch$ and let $Q$ be the cube associated with a Zariski cover 
$\{U_i\to X\}_{1\leq i\leq n}$ in $\lSm/S$.
Then, in $\infShv_{sNis}(\lSm/S)[\boxx^{-1}]$, we have an equivalence
\[
\colimit(\horn(Q))
\simeq
X.
\]
\end{prop}

\begin{equation}
\label{Hoch.37.1}
\TC^{-}(A)
:=
\THH(A)^{h\mathbb{T}},
\text{ }
\TP(A)
:=
\THH(A)^{t\mathbb{T}},
\text{ }
\TC(A)
:=
\TC(\THH(A)).
\end{equation}

\begin{prop}
\label{Hoch.9}
There are $\mathbb{T}$-equivariant equivalences
\[
\Bcy (\N,j)
\simeq
\left\{
\begin{array}{ll}
* & \text{if }j=0
\\
\mathbb{T}/C_j & \text{if }j>1
\end{array}
\right.
\]
and
\[
\Bcy (\Z,j)
\simeq
\mathbb{T}/C_j
\;
(j\in \Z),
\;
\Brep(\N,j)
\simeq
\mathbb{T}/C_j
\;
(j\geq 0),
\]
where $C_0$ denotes the trivial group for convenience.
The naturally induced maps
\[
\Bcy (\N,j)
\to
\Bcy(\Z,j)
\;
(j\geq 1)
\text{ and }
\Brep (\N,j)
\to
\Bcy(\Z,j)
\;
(j\geq 0)
\]
are $\mathbb{T}$-equivariant equivalences.
\end{prop}

\begin{lem}
\label{Hoch.2}
For all monoids $M$ and $N$, there is a canonical $\mathbb{T}$-equivariant equivalence
\begin{equation}
\label{Hoch.2.1}
\Bcy (M\times N) \cong \Bcy M \times \Bcy N.
\end{equation}
\end{lem}

\begin{lem}
\label{Hoch.14}
For $Q:=\N^r\oplus \Z^s$, $P:=\N^r \oplus \N \oplus \Z^{s-1}$, and $M:=0\oplus \N \oplus \Z^{s-1}$, 
there is a canonical $\mathbb{T}$-equivariant equivalence
\[
\Brep(P,M)
\simeq
\Bcy Q\times_Q P.
\]
\end{lem}

\begin{prop}
\label{Hoch.12}
Let $(A,M)$ be an integral pre-log ring.
Then the naturally induced map of cyclotomic spectra
\begin{equation}
\label{Hoch.12.2}
\logTHH(A,M)
\to
\logTHH(A,M^a)
\end{equation}
is an equivalence.
\end{prop}

\begin{prop}
\label{Hoch.11}
Let $(R,M)$ be a pre-log ring, and let $N\to P$ be a homomorphism of monoids.
Then there is a canonical equivalence of $\E_\infty$-rings in cyclotomic spectra
\begin{equation}
\label{Hoch.11.1}
\logTHH((R,M)\otimes (\Z[P],N))
\simeq
\logTHH(R,M)
\otimes_{\Striv}
\Sphere[\Brep(P,N)].
\end{equation}
\end{prop}

\begin{equation}
\label{Thomdf.5.3}
\tcofib (Q|_{(\Delta^1)^{i-1}\times \{0\}\times (\Delta^1)^{n-i}})
\to
\tcofib (Q|_{(\Delta^1)^{i-1}\times \{1\}\times (\Delta^1)^{n-i}})
\to
\tcofib(Q).
\end{equation}

\begin{prop}
\label{Thomdf.6}
Let $i_1\colon X_{01}\to X_{11}$, $\ldots$, $i_n\colon X_{0n}\to X_{1n}$ be maps in a pointed symmetric monoidal 
$\infty$-category $\cC$ with small colimits.
If the tensor product operation on $\cC$ preserves small colimits in each variable, 
then there is a canonical equivalence in $\cC$
\[
\cofib(i_1)\otimes \cdots \otimes \cofib(i_n)
\simeq
\tcofib(Q_{i_1\cdots i_n}).
\]
Here we write $Q$ for the naturally associated $n$-cube sending any point $(a_1,\ldots,a_n)$ to 
$X_{a_1 1}\otimes \cdots \otimes X_{a_n n}$.
\end{prop}

\begin{prop}
\label{Thomdf.15}
Let $\cC$ be a stable $\infty$-category.
Then an $n$-cube $Q$ in $\cC$ is cartesian if and only if it is cocartesian.
\end{prop}

\begin{thm}
\label{Hoch.13}
If $X\in \lSch$, then the naturally induced morphism of spectra
\[
\logTHH(X)
\to
\logTHH(X\times (\P^n,\P^{n-1}))
\]
is an equivalence for all $n\geq 1$.
The same holds for $\logTC^-$, $\logTP$, and $\logTC$.
\end{thm}

\begin{proof}
Owing to Propositions \ref{Hoch.12} and \ref{Hoch.11}, we are reduced to consider $X=\Spec{\Z}$.
The fs log scheme $(\P^n,\P^{n-1})$ admits the Zariski cover $\{U_0,\ldots,U_n\}$,
where
\[
U_0:=\Spec{\Z[x_1/x_0,\ldots,x_n/x_0]}
\]
and
\[
U_i:=\Spec{\Z[x_0/x_i,\ldots,x_n/x_i],\N(x_0/x_i)}
\]
for $1\leq i\leq n$.
For every nonempty subset $I=\{i_1,\ldots,i_r\}\subset \{0,\ldots,n\}$, we set
\[
U_I:=U_{i_1}\cap \cdots \cap U_{i_r}.
\]
Let $\{e_0,\ldots,e_n\}$ be the standard basis in $\Z^{n+1}$, 
and let $P_I$ be the submonoid of $\Z^{n+1}$ generated by $e_i-e_j$ for all $i\in \{0,\ldots,n\}$ and $j\in I$.
We set $P_\emptyset := 0$ and let 
$F_I$ be the face of $P_I$ generated by $e_0-e_j$ for all $j\in I$ if $0\in I$.
By convention, 
$F_I=0$ if $0\notin I$.
There is an isomorphism of fs log schemes
\[
U_I
\cong
\Spec{F_I\to \Z[P_I]}
\]
if $I\neq \emptyset$, 
and we also have
\[
\Spec{\Z}
\cong
\Spec{F_\emptyset \to \Z[P_\emptyset]}.
\]
Let $C_n$ be the $(n+1)$-cube in $\infSpc$ naturally defined by $\Brep(P_I,F_I)$ for all $I\subset \{0,\ldots,n\}$,
which naturally induces an $(n+1)$-cube $\Sphere[C_n]$ in $\infSpt$.
We need to show that there is an equivalence of spectra
\[
\tfib(\Sphere[C_n])
\simeq
0,
\]
which implies the claim by Proposition \ref{logHprop.3}.
Owing to Proposition \ref{Thomdf.15}, it suffices to show there is an equivalence of spectra
\[
\tcofib(\Sphere[C_n])
\simeq
0.
\]
Since $\Sphere[-]$ is a left adjoint,
it preserves small colimits.
Hence it suffices to show there is an equivalence of spaces
\begin{equation}
\label{Hoch.13.6}
\tcofib(C_n)
\simeq
0.
\end{equation}

\

It is convenient to consider $Q_J:=P_{\{0,\ldots,n\}-J}$ and $G_J:=F_{\{0,\ldots,n\}-J}$ for subsets 
$J\subset \{0,\ldots,n\}$.
We set $Q_j:=Q_{\{j\}}$ for $j=0,\ldots,n$.
Let $\{f_1,\ldots,f_n\}$ be the standard basis in $\Z^n$.
By identifying $e_j-e_0$ with $f_j$ for $j=1,\ldots,n$, we have isomorphisms of monoids
\[
Q_j
\cong
\{ (x_1,\ldots,x_n)\in \Z^n : x_j\geq 0 \}
\]
for $j=1,\ldots,n$ and
\[
Q_0
\cong
\{ (x_1,\ldots,x_n)\in \Z^n : x_1+\cdots+x_n \leq 0\}.
\]
Furthermore, we have an isomorphism of monoids
\[
Q_J
\cong
\bigcap_{j\in J} Q_j.
\]

Suppose $0\notin J$.
We set $M_{j,J}:=\N f_j$ if $j\in J$ and $M_{j,J}:=\Z f_j$ if $j\notin J$.
Then there is a canonical isomorphism of monoids
\[
Q_J\cong M_{1,J}\times \cdots \times M_{n,J}.
\]
Furthermore, we have $G_J=0$.
Lemma \ref{Hoch.2} gives a canonical equivalence of spaces
\begin{equation}
\label{Hoch.13.1}
\Brep (Q_J,G_J)
\simeq
\Bcy M_{1,J}
\times
\cdots
\times
\Bcy M_{n,J}.
\end{equation}

On the other hand, suppose $0\in J$.
Let $i$ be any element in $\{0,\ldots,n\}-J$.
We set $M_{j,J}':=\N(f_j-f_i)$ if $j\in J$ and $M_{j,J}':=\Z(f_j-f_i)$ if $j\notin J$.
Then there are isomorphisms of monoids
\[
Q_J
\cong
M_{0,J}'\times M_{1,J}'\times \cdots \times M_{i-1,J}'\times M_{i+1,J}'\times \cdots \times M_{n,J}'
\]
and
\[
Q_{J-\{0\}}
\cong
\Z\times M_{1,J}'\times \cdots \times M_{i-1,J}'\times M_{i+1,J}'\times \cdots \times M_{n,J}'.
\]
Furthermore, $G_J$ is the face of $Q_J$ generated by $f_0-f_i\in M_{0,J}'$.
From Lemma \ref{Hoch.14}, we have a canonical equivalence of spaces
\begin{equation}
\label{Hoch.13.3}
\Brep (Q_J,G_J)
\simeq
\Brep (Q_{J-\{0\}},G_{J-\{0\}})\times_{Q_{J-\{0\}}}Q_J.
\end{equation}
Together with $Q_J\cong Q_{J-\{0\}}\times_{\Z^n} Q_0$, we have a canonical equivalence of spaces
\begin{equation}
\label{Hoch.13.5}
\Brep (Q_J,G_J)
\simeq
\Brep (Q_{J-\{0\}},G_{J-\{0\}}) \times_{\Z^n}Q_0.
\end{equation}

\

Let $V$ be the cofiber of $\Bcy \N\to \Bcy \Z$.
Owing to Proposition \ref{Hoch.9}, there is a decomposition
\[
V
\simeq
\coprod_{j\leq 0} V_j
\]
with $V_0:=\cofib(*\to S_1/C_0)$ and $V_j:=S_1/C_j$ for $j<0$.
In particular, the canonical map $V\to \Z$ factors through $-\N\subset \Z$.
Let $D_n$ (resp.\ $D_n'$) be the $n$-cube naturally associated with $\Brep (Q_J,G_J)$ for all $J$ with $0\notin J$ (resp.\ $0\in J$).
From \eqref{Hoch.13.1}, we see that $D_n$ is the $n$-cube formulated in Proposition \ref{Thomdf.6} associated with the maps
\[
\Bcy (\N f_i)
\to
\Bcy (\Z f_i)
\]
for $i=1,\ldots,n$.
Proposition \ref{Thomdf.6} gives an equivalence of spaces
\[
\tcofib D_n
\simeq
V\times \cdots \times V
\;
(n\text{ copies of }V).
\]
Together with \eqref{Hoch.13.5} we have an equivalence of spaces
\[
\tcofib D_n'
\simeq
(V\times \cdots \times V)\times_{\Z^n}Q_0
\;
(n\text{ copies of }V).
\]
The map $(V\times \cdots \times V)\to \Z^n$ factors through $(-\N)^n$, and we have $(-\N)^n\subset Q_0$.
Hence we have an equivalence of spaces
\begin{equation}
\label{Hoch.13.4}
\tcofib D_n'
\simeq
\tcofib D_n.
\end{equation}
From \eqref{Thomdf.5.3}, we have a cofiber sequence
\[
\tcofib D_n\to \tcofib D_n'\to \tcofib C_n.
\]
Together with \eqref{Hoch.13.4} we deduce $\tcofib C_n\simeq 0$.

\

Finally, we use \eqref{Hoch.37.1} to deduce the claims for $\logTC$ and $\logTC^-$.
\end{proof}

\end{document}
